\section{RLP 的动机：从观察文本到通过思考学习}

RLP 不再只调整数据顺序，而是在每个选定 token 前引入一个可采样的 thought action。核心思想是：模型先产生若干内部 reasoning trace，再预测真实下一 token；如果某条 thought 提高了该 token 的概率，就获得正奖励。这样，普通预训练文本同时提供 target token 与自监督 reward reference。

\subsection{The Tale of Two Learners}

Leo 先搭一个粗糙但可通过小车的桥，再从反馈中改进；Bolt 阅读所有桥梁资料，造出漂亮却不满足约束的模型。故事强调 learning by doing 与 learning by observing 的差别。Next-token prediction 擅长从已写文本中模仿统计模式，RLP 则试图给模型一次“先提出解释、再用真实 token 检验”的探索机会。

\lecturefigure{slide-023.jpg}{两个学习者：动手试错与只观察范例的差别}{Stanford 官方录像，00:28:14--00:29:37}

\begin{warningbox}{Thought 不是自动可信的解释}
模型生成的 chain-of-thought 可能包含事后合理化、错误事实或不可解释捷径。RLP 只检查 thought 是否提高 observed next token 的预测概率，不验证它是否符合人类推理、是否因果忠实、是否能在分布外成立。
\end{warningbox}

\subsection{标准预训练的问题不是没有梯度，而是没有探索动作}

标准 NTP 已经提供密集梯度，但 action space 只有对下一 token 分布的直接调整；复杂 reasoning 通常等到 SFT、RLHF 或 RLVR 才显式出现。RLP 的问题陈述是：能否把 exploration 前移，同时保持普通文档流的可扩展性？这一动机也回应数据墙——如果新 token 越来越贵，就要提高每个 token 产生的学习信号。

\lecturefigure{slide-024.jpg}{标准预训练把显式 exploration 推迟到 post-training}{Stanford 官方录像，00:29:37--00:31:23}

标准 NTP 目标为
\[
\mathcal{L}_{\mathrm{NTP}}(\theta)
=-\sum_{t=1}^{L}\log p_\theta(x_t\mid x_{<t}).
\]
其中，$x_{<t}$ 表示上下文，$x_t$ 是语料中的真实下一 token，$p_\theta$ 是模型分布。该目标会学到大量结构，但没有单独变量表示“在输出前探索了哪条 thought”。

\subsection{Vanilla 与 RLP 的条件概率差别}

在 photosynthesis 例子里，vanilla model 直接从上下文预测 `sunlight`；RLP model 先生成“该过程依赖 solar energy”的 thought，再条件于 context 与 thought 预测同一个 token。关键不是多输出一段文字，而是把 thought 视为可优化 action，并比较它相对 context-only baseline 的增益。

\lecturefigure{slide-025.jpg}{Vanilla next-token prediction 与 thought-conditioned RLP prediction}{Stanford 官方录像，00:31:23--00:32:34}

\[
p_\theta(x_t\mid x_{<t},c_t)
\quad\text{versus}\quad
p_\phi(x_t\mid x_{<t}),
\]
其中，$c_t$ 是在位置 $t$ 前采样的 thought，$\theta$ 是 thought policy / reasoned predictor 参数，$\phi$ 是 no-think baseline 参数。两者看到相同 context 和真实 target，差别只在是否加入 $c_t$。

\begin{importantbox}{RLP 的自监督闭环}
语料中的真实 token $x_t$ 同时承担两项角色：它是 NTP target，也是评估 thought 是否有 predictive utility 的参照。RLP 因此不需要人工答案 verifier，但仍依赖真实下一 token，并不是无监督地奖励任何“看起来像推理”的文本。
\end{importantbox}

\subsection{本章小结}

RLP 把 thought 变成预测前的 exploratory action，用真实下一 token 检查该 action 是否提供信息。它与 NTP 的差别不是有没有训练信号，而是有没有显式探索路径和对路径质量的相对奖励。下面逐步拆解完整训练环。

